\chapter{How to read this, and what the project was}

\section{The one-paragraph version}

A group at MIT trained a large language model --- the same kind of machine that
powers a chatbot, but fed protein sequences instead of English --- to do two
things with spider silk proteins. Given a protein sequence, predict how strong
and stretchy the resulting silk fibre will be. And, in reverse, given a wish
list of mechanical properties, invent a protein sequence that should have them.
They published numbers showing the two directions agree with each other. I
downloaded their model, rebuilt their dataset from raw sources, re-ran their
procedure, and then ran four further experiments they did not run, to work out
what their headline numbers actually measure.

\section{Why reproduce anything at all}

A published result is a claim about what would happen if someone else did the
same thing. Most of the time nobody checks. When somebody does check, three
outcomes are possible, and all three are useful:

\begin{enumerate}[leftmargin=*,itemsep=3pt]
  \item \term{It reproduces.} The claim is now stronger than it was, because it
        has survived an attempt to break it.
  \item \term{It reproduces, but the number means something narrower than it
        appeared to.} This is the most common outcome in machine learning, and
        it is what happened here. Nothing was wrong; the number was
        under-specified.
  \item \term{It does not reproduce.} Rare, and worth knowing.
\end{enumerate}

This project landed mostly in category~1 with an important slice of
category~2. The model reproduces \emph{exactly} --- to four decimal places, on
the paper's own inputs. What does not reproduce is my own attempt to
\emph{search} for good designs, and understanding why that is a different thing
from the model failing took most of the work.

\section{The shape of the argument}

Figure~\ref{fig:roadmap} is the whole document on one page. Nothing in it will
mean much yet; come back to it after Chapter~\ref{ch:paper}.

\begin{figure}[htbp]
\centering
\begin{tikzpicture}[
  font=\small,
  box/.style={draw=greyc!70,fill=boxbg,rounded corners=2pt,inner sep=6pt,
              text width=3.1cm,align=center,minimum height=1.0cm},
  pbox/.style={box,draw=papercol,fill=papercol!8},
  obox/.style={box,draw=ourscol,fill=ourscol!8},
  ar/.style={-{Stealth[length=2.5mm]},thick,greyc!80},
]
\node[pbox] (paper) {\textbf{The paper}\\ trains a model, reports \Rsq};
\node[obox,below=0.9cm of paper] (rebuild) {\textbf{Rebuild}\\ dataset, constants, checkpoint};
\node[obox,below=0.9cm of rebuild] (repro) {\textbf{Re-run}\\ the design procedure};
\node[obox,right=1.5cm of repro] (short) {\textbf{Shortfall}\\ on all 8 targets};
\node[obox,right=1.5cm of short] (exact) {\textbf{Exact match}\\ on the paper's own sequences};
\node[obox,below=0.9cm of repro] (ext) {\textbf{Four extensions}\\ what is the model doing?};
\node[obox,below=0.9cm of ext] (mem) {\textbf{Answer}\\ largely retrieval};

\draw[ar] (paper) -- (rebuild);
\draw[ar] (rebuild) -- (repro);
\draw[ar] (repro) -- (short);
\draw[ar] (short) -- (exact);
\draw[ar] (repro) -- (ext);
\draw[ar] (ext) -- (mem);

\node[right=0.25cm of exact,text width=3.0cm,align=left,font=\scriptsize\itshape,
      text=ourscol]
  {so the gap is in the \emph{search}, not the model};
\node[left=0.25cm of mem,text width=2.6cm,align=right,font=\scriptsize\itshape,
      text=ourscol]
  {71\% of training labels returned verbatim};
\end{tikzpicture}
\caption[Roadmap]{The argument in one figure. Blue is the original paper;
orange is this project. Read it again after Chapter~\ref{ch:paper}.}
\label{fig:roadmap}
\end{figure}

\section{Conventions used throughout}

\begin{itemize}[leftmargin=*,itemsep=2pt]
  \item \paperterm{Blue bold} marks something belonging to the original paper.
  \item \ourterm{Orange bold} marks something this project added or decided.
  \item Amino-acid sequences are set as \aaseq{GAGAGSGAAG}.
  \item Every number I quote as a result was produced by a script in the
        project and is traceable to a file; the mapping is in
        Chapter~\ref{ch:reference}.
\end{itemize}

Three kinds of highlighted box appear:

\begin{keybox}{A concept box looks like this}
Used for definitions and for ideas that need to be held in mind for several
pages.
\end{keybox}

\begin{warnbox}{A caution box looks like this}
Used for traps, for things I got wrong, and for places where a number is easy
to over-read.
\end{warnbox}

\begin{goodbox}{A result box looks like this}
Used for measured outcomes.
\end{goodbox}

\section{A note on the three retractions}

Three times in this project I stated something, and it was wrong, and I had to
withdraw it. They are not hidden in a footnote --- Chapter~\ref{ch:mistakes} is
about them, because they turned out to be the most transferable thing I
learned. Briefly, they were:

\begin{enumerate}[leftmargin=*,itemsep=3pt]
  \item A claim that three of the paper's motif definitions contradicted the
        source they cite. \term{Cause:} I read a two-column PDF table as
        linear text, and the columns interleaved.
  \item A set of \Rsq{} values suggesting the paper's model failed on the
        paper's own sequences. \term{Cause:} my extraction of those sequences
        was silently corrupt, in a way that made every length come out right.
  \item A finding that a composition analysis recovered 12 of 14 published
        residues. \term{Cause:} the selection rule could not fail; random
        noise scored better.
\end{enumerate}

The pattern is the useful part. All three survived internal consistency
checks and all three were caught by comparison against an \emph{external}
published quantity that had played no part in generating the claim.
